% GENERATED FILE. Edit canonical lessons, then run npm run generate:books.
% canonical-source-hash: fnv1a64:a5c5a6af454613b6
% canonical-lessons: MW-C102-sagla, MW-C102-dujo, MW-C102-apno, MW-C102-koi, MW-C102-harek

\chapter{All, Another, Our own, Someone, Each}
\label{ch:mw-all-another-our-own-someone-each}
\hlchaptermodality{102}

\begin{chapteropening}
\textbf{By the end of this chapter:} \emph{I can say all, another, our own, someone and each.}

Complete the last lesson of chapter 102: I can say all, another, our own, someone and each.
\end{chapteropening}

\section[\texorpdfstring{\mw{सगळा} (sagḷā)}{सगळा (sagḷā)}]{\mw{सगळा} (sagḷā) — all, everyone}

\label{lesson:MW-C102-sagla}



\begin{quote}
Before the new one: say the Marwadi for which, then the Marwadi for who.
\end{quote}

\subsection*{You'll want to know: \mw{सगळा}}
\textbf{\mw{सगळा}} — \emph{sagḷā} — \textquotedblleft{}all, everyone\textquotedblright{}.

\begin{cousinweb}
\textbf{\mw{सगळा}} (\emph{sagḷā}) is usually linked with Sanskrit \emph{sakala}, \textquotedblleft{}whole\textquotedblright{}, under which CDIAL lists Old Gujarati \emph{sayala}, \textquotedblleft{}all\textquotedblright{} (CDIAL 13066). CDIAL gives no Marwari form, and that entry does not explain the \emph{g} and \emph{ḷ} of \textbf{\mw{सगळा}} (\emph{sagḷā}), so the link is not settled.
\end{cousinweb}

\begin{grammarlens}[title={What you've built}]
One more word for this chapter.
\end{grammarlens}

\subsection*{Your turn}
\begin{itemize}
\raggedright
  \item \emph{Say it:} \emph{sagḷā}
  \item \emph{Say it:} \emph{sagḷā}, once more
  \item \emph{Say it:} say \emph{sagḷā}
  \item \emph{Recall:} say the Marwadi for which, then the Marwadi for who, then say \emph{sagḷā} again
\end{itemize}

\begin{tcolorbox}[breakable,colback=teal!4,colframe=teal!35!black,title={Before you move on}]
What does \mw{सगळा} mean? (\textquotedblleft{}All, everyone\textquotedblright{}.) Say it once more.
\end{tcolorbox}
\section[\texorpdfstring{\mw{दूजो} (dūjo)}{दूजो (dūjo)}]{\mw{दूजो} (dūjo) — another, the second}

\label{lesson:MW-C102-dujo}



\begin{quote}
Before the new one: say the Marwadi for who, then the Marwadi for all.
\end{quote}

\subsection*{You'll want to know: \mw{दूजो}}
\textbf{\mw{दूजो}} — \emph{dūjo} — \textquotedblleft{}another, the second\textquotedblright{}.

\begin{cousinweb}
\textbf{\mw{दूजो}} (\emph{dūjo}) is a tadbhava, a word inherited and reshaped by sound change. It comes from \emph{*dutīya}, \textquotedblleft{}second\textquotedblright{}, a form CDIAL reconstructs from its descendants. CDIAL gives no Marwari form in this entry (CDIAL 6402); the match is through Hindi \emph{dūjā} and Gujarati \emph{dujũ}. CDIAL sets it beside Sanskrit \emph{dvitīya}, \textquotedblleft{}second\textquotedblright{}.
\end{cousinweb}

\begin{grammarlens}[title={What you've built}]
One more word for this chapter.
\end{grammarlens}

\subsection*{Your turn}
\begin{itemize}
\raggedright
  \item \emph{Say it:} \emph{dūjo}
  \item \emph{Say it:} \emph{dūjo}, once more
  \item \emph{Say it:} say \emph{dūjo}
  \item \emph{Recall:} say the Marwadi for who, then the Marwadi for all, then say \emph{dūjo} again
\end{itemize}

\begin{tcolorbox}[breakable,colback=teal!4,colframe=teal!35!black,title={Before you move on}]
What does \mw{दूजो} mean? (\textquotedblleft{}Another, the second\textquotedblright{}.) Say it once more.
\end{tcolorbox}
\section[\texorpdfstring{\mw{आपणो} (āpṇo)}{आपणो (āpṇo)}]{\mw{आपणो} (āpṇo) — our own, one's own}

\label{lesson:MW-C102-apno}



\begin{quote}
Before the new one: say the Marwadi for all, then the Marwadi for another.
\end{quote}

\subsection*{You'll want to know: \mw{आपणो}}
\textbf{\mw{आपणो}} — \emph{āpṇo} — \textquotedblleft{}our own, one's own\textquotedblright{}.

\begin{cousinweb}
\textbf{\mw{आपणो}} (\emph{āpṇo}) is a tadbhava, a word inherited and reshaped by sound change. It comes from \emph{*ātmanaka}, \textquotedblleft{}own\textquotedblright{}, a form CDIAL reconstructs from its descendants. CDIAL lists Old Marwari \emph{āpaṇau} under it (CDIAL 1135). CDIAL reconstructs \emph{*ātmanaka} from Sanskrit \emph{ātman}, \textquotedblleft{}self\textquotedblright{}.
\end{cousinweb}

\begin{grammarlens}[title={What you've built}]
One more word for this chapter.
\end{grammarlens}

\subsection*{Your turn}
\begin{itemize}
\raggedright
  \item \emph{Say it:} \emph{āpṇo}
  \item \emph{Say it:} \emph{āpṇo}, once more
  \item \emph{Say it:} say \emph{āpṇo}
  \item \emph{Recall:} say the Marwadi for all, then the Marwadi for another, then say \emph{āpṇo} again
\end{itemize}

\begin{tcolorbox}[breakable,colback=teal!4,colframe=teal!35!black,title={Before you move on}]
What does \mw{आपणो} mean? (\textquotedblleft{}Our own, one's own\textquotedblright{}.) Say it once more.
\end{tcolorbox}
\section[\texorpdfstring{\mw{कोई} (koī)}{कोई (koī)}]{\mw{कोई} (koī) — someone, any}

\label{lesson:MW-C102-koi}



\begin{quote}
Before the new one: say the Marwadi for another, then the Marwadi for our own.
\end{quote}

\subsection*{You'll want to know: \mw{कोई}}
\textbf{\mw{कोई}} — \emph{koī} — \textquotedblleft{}someone, any\textquotedblright{}.

\begin{cousinweb}
\textbf{\mw{कोई}} (\emph{koī}) is a tadbhava, a word inherited from Sanskrit and reshaped by sound change. It comes from \emph{kaścid}, \textquotedblleft{}some one, anyone\textquotedblright{}. CDIAL lists Old Marwari \emph{koi} under it (CDIAL 2967). CDIAL notes that it partly ran together with \emph{ko'pi}, \textquotedblleft{}whoever\textquotedblright{}.
\end{cousinweb}

\begin{grammarlens}[title={What you've built}]
One more word for this chapter.
\end{grammarlens}

\subsection*{Your turn}
\begin{itemize}
\raggedright
  \item \emph{Say it:} \emph{koī}
  \item \emph{Say it:} \emph{koī}, once more
  \item \emph{Say it:} say \emph{koī}
  \item \emph{Recall:} say the Marwadi for another, then the Marwadi for our own, then say \emph{koī} again
\end{itemize}

\begin{tcolorbox}[breakable,colback=teal!4,colframe=teal!35!black,title={Before you move on}]
What does \mw{कोई} mean? (\textquotedblleft{}Someone, any\textquotedblright{}.) Say it once more.
\end{tcolorbox}
\section[\texorpdfstring{\mw{हरेक} (harek)}{हरेक (harek)}]{\mw{हरेक} (harek) — each, every}

\label{lesson:MW-C102-harek}



\begin{quote}
Before the new one: say the Marwadi for our own, then the Marwadi for someone.
\end{quote}

\subsection*{You'll want to know: \mw{हरेक}}
\textbf{\mw{हरेक}} — \emph{harek} — \textquotedblleft{}each, every\textquotedblright{}.

\begin{cousinweb}
\textbf{\mw{हरेक}} (\emph{harek}) joins \emph{har}, \textquotedblleft{}every\textquotedblright{}, usually traced to Persian \emph{har}, with \emph{ek}, \textquotedblleft{}one\textquotedblright{}, from Sanskrit \emph{eka} (CDIAL 2462).
\end{cousinweb}

\begin{grammarlens}[title={What you've built}]
One more word for this chapter.
\end{grammarlens}

\subsection*{Your turn}
\begin{itemize}
\raggedright
  \item \emph{Say it:} \emph{harek}
  \item \emph{Say it:} \emph{harek}, once more
  \item \emph{Say it:} say \emph{harek}
  \item \emph{Recall:} say the Marwadi for our own, then the Marwadi for someone, then say \emph{harek} again
\end{itemize}

\begin{tcolorbox}[breakable,colback=teal!4,colframe=teal!35!black,title={Before you move on}]
What does \mw{हरेक} mean? (\textquotedblleft{}Each, every\textquotedblright{}.) Say it once more.
\end{tcolorbox}
